% =====================================================================
%  Phương pháp T5 sinh cho Trích xuất Aspect Term (ATE)
%  Dùng % =====================================================================
%  Phương pháp T5 sinh cho Trích xuất Aspect Term (ATE)
%  Dùng % =====================================================================
%  Phương pháp T5 sinh cho Trích xuất Aspect Term (ATE)
%  Dùng % =====================================================================
%  Phương pháp T5 sinh cho Trích xuất Aspect Term (ATE)
%  Dùng \input{phuong_phap_ate} trong main.tex
% =====================================================================

% =====================================================================
\section{Phương pháp T5 sinh cho Trích xuất Aspect Term}
\label{sec:ate}
% =====================================================================

Luận văn sử dụng hai cách tiếp cận dựa trên T5 để trích xuất aspect
term (ATE), đều được đánh giá theo cùng một độ đo ATE~F1:

\begin{enumerate}
  \item \textbf{GAS} (\textit{Generative Aspect Sentiment}) —
    mô hình seq2seq sinh đồng thời ba nhãn
    $(a_i, c_i, s_i)$ trong một lần giải mã; khi dùng cho ATE, chỉ
    trường $a_i$ được giữ lại
    (Mục~\ref{subsec:gas_full}).

  \item \textbf{T5-ATE} — mô hình seq2seq huấn luyện chuyên biệt cho
    ATE: chuỗi đích chỉ chứa aspect term, không có category hay
    sentiment (Mục~\ref{subsec:t5ate}).
\end{enumerate}

% =====================================================================
\subsection{Mô hình GAS (Generative Aspect Sentiment)}
\label{subsec:gas_full}
% =====================================================================

GAS~\cite{zhang2021gas} xây dựng bài toán ABSA như một bài toán dịch
chuỗi, với mô hình nền là
\texttt{T5ForConditionalGeneration}~\cite{raffel2020t5} phiên bản
\texttt{t5-base}.  Với mỗi câu đầu vào $\mathbf{x}$, mô hình được
huấn luyện để sinh ra chuỗi đích $\mathbf{y}$ theo cú pháp cố định:

\begin{equation}
  \mathbf{y} = \texttt{(}a_1,\, c_1,\, s_1\texttt{);\;
               (}a_2,\, c_2,\, s_2\texttt{);\;} \dots
  \label{eq:gas_format}
\end{equation}

\noindent trong đó $a_i$ là cụm từ khía cạnh (\textit{aspect term}),
$c_i \in \{\text{SERVICE, FACILITY, AMENITY, EXPERIENCE, BRANDING,
LOYALTY}\}$ là danh mục khía cạnh, và
$s_i \in \{\text{positive, negative, neutral}\}$ là cực tính cảm xúc.
Nếu câu không chứa khía cạnh nào, mô hình sinh ra chuỗi
\texttt{"none"}.

% ---------------------------------------------------------------------
\subsubsection{Tiền xử lý dữ liệu}
\label{subsubsec:gas_data}
% ---------------------------------------------------------------------

Dữ liệu đầu vào được đọc từ định dạng \texttt{.apc} bốn dòng lặp lại
(câu — aspect term — aspect category — sentiment).
Vì định dạng này lưu một dòng cho mỗi cặp \textit{(câu, aspect)},
các mẫu cùng thuộc một câu được nhóm lại thành một
\textit{record} duy nhất.  Chuỗi đích tương ứng được xây dựng bằng
cách nối tất cả bộ ba của câu đó theo cú pháp~\eqref{eq:gas_format}.

Mỗi record có dạng:

\begin{equation}
  r = \bigl(\underbrace{x}_{\text{input\_text}},\;
            \underbrace{y}_{\text{target\_text}},\;
            \underbrace{\{a_i\}}_{\text{aspects}},\;
            \underbrace{\{(a_i,\,c_i,\,s_i)\}}_{\text{gold\_triples}}
      \bigr)
  \label{eq:gas_record}
\end{equation}

Tokenization được thực hiện bởi \texttt{T5Tokenizer} với
\texttt{max\_length}$= 128$ cho cả đầu vào lẫn đầu ra.  Các vị trí
padding trong chuỗi nhãn được thay bằng $-100$ để hàm mất mát
cross-entropy bỏ qua.

% ---------------------------------------------------------------------
\subsubsection{Huấn luyện}
\label{subsubsec:gas_train}
% ---------------------------------------------------------------------

Mô hình được tinh chỉnh bằng hàm mất mát cross-entropy seq2seq chuẩn
(T5 tự tính theo cơ chế \textit{teacher forcing}):

\begin{equation}
  \mathcal{L}_{\text{GAS}} =
    -\sum_{t=1}^{T} \log P\!\left(y_t \mid y_{<t},\, \mathbf{x};\,
    \Theta\right)
  \label{eq:gas_loss}
\end{equation}

\noindent với $T$ là độ dài chuỗi đích và $\Theta$ là toàn bộ tham số
mô hình.  Quá trình tối ưu sử dụng AdamW với learning rate
$3 \times 10^{-4}$, kết hợp \textit{linear warmup} trong 10\% tổng
số bước.  Kỹ thuật mixed precision (AMP) được bật khi có GPU CUDA.

Early stopping theo dõi \textbf{joint F1} trên tập phát triển
(patience$= 4$ epoch): một dự đoán được tính là True Positive khi
\textit{cả ba} trường $(a_i, c_i, s_i)$ đều khớp chính xác với nhãn
vàng tương ứng.

% ---------------------------------------------------------------------
\subsubsection{Hậu xử lý và trích xuất aspect term}
\label{subsubsec:gas_postprocess}
% ---------------------------------------------------------------------

Chuỗi văn bản do T5 sinh ra được phân tích cú pháp bằng biểu thức
chính quy ba nhóm:

\begin{equation}
  \texttt{\_GAS\_RE} =
    \texttt{\textbackslash(\textbackslash s*(.+?),\textbackslash s*%
    (.+?),\textbackslash s*(positive|negative|neutral)\textbackslash s*%
    \textbackslash)}
  \label{eq:gas_regex}
\end{equation}

\noindent Group~1, 2, 3 lần lượt trả về $a_i$, $c_i$, $s_i$.  Để
sử dụng GAS như một baseline ATE, chỉ Group~1 ($a_i$) được giữ lại;
$c_i$ và $s_i$ bị bỏ qua ở bước đánh giá.  Nếu chuỗi sinh là
\texttt{"none"} hoặc không khớp mẫu nào, hàm trả về danh sách rỗng.

Sau khi parse, bước chuẩn hóa n-gram Levenshtein được áp dụng cho
$a_i$ (mục~\ref{subsubsec:levenshtein}); $c_i$ và $s_i$ giữ nguyên.

% =====================================================================
\subsection{Mô hình T5-ATE}
\label{subsec:t5ate}
% =====================================================================

Khác với GAS, T5-ATE được huấn luyện chuyên biệt cho ATE: chuỗi đích
chỉ chứa aspect term, không có category hay sentiment.  Điều này rút
ngắn không gian sinh và giảm áp lực cho decoder khi chỉ cần học một
nhãn.

% ---------------------------------------------------------------------
\subsubsection{Kiến trúc và định dạng đầu ra}
\label{subsubsec:t5ate_arch}
% ---------------------------------------------------------------------

Mô hình nền cũng là \texttt{T5ForConditionalGeneration} phiên bản
\texttt{t5-base}.  Với mỗi câu đầu vào $\mathbf{x}$, mô hình sinh ra
chuỗi đích chỉ chứa aspect term:

\begin{equation}
  \mathbf{y} = \texttt{(}a_1\texttt{);\;(}a_2\texttt{);\;} \dots
  \label{eq:ate_format}
\end{equation}

\noindent Nếu câu không chứa khía cạnh nào, mô hình sinh ra chuỗi
\texttt{"none"}.  Quá trình giải mã sử dụng \textit{beam search} với
chiều rộng tia $k = 4$.

% ---------------------------------------------------------------------
\subsubsection{Tiền xử lý dữ liệu}
\label{subsubsec:t5ate_data}
% ---------------------------------------------------------------------

Dữ liệu được đọc từ cùng định dạng \texttt{.apc} nhưng các trường
category và sentiment bị bỏ qua hoàn toàn.  Các mẫu cùng thuộc một
câu được nhóm lại; chuỗi đích được xây dựng theo
cú pháp~\eqref{eq:ate_format}.  Mỗi record có dạng:

\begin{equation}
  r = \bigl(\underbrace{x}_{\text{input\_text}},\;
            \underbrace{y}_{\text{target\_text}},\;
            \underbrace{\{a_i\}}_{\text{aspects}}
      \bigr)
  \label{eq:ate_record}
\end{equation}

Tokenization sử dụng \texttt{max\_length}$= 128$ cho đầu vào và
$= 64$ cho đầu ra — ngắn hơn GAS vì không cần mã hóa category và
sentiment.

% ---------------------------------------------------------------------
\subsubsection{Huấn luyện}
\label{subsubsec:t5ate_train}
% ---------------------------------------------------------------------

Hàm mất mát và thiết lập tối ưu giống GAS~\eqref{eq:gas_loss}:
AdamW với learning rate $3 \times 10^{-4}$ và linear warmup 10\%.
Điểm khác biệt: T5-ATE \textbf{không áp dụng early stopping} — mô
hình được huấn luyện đủ 20 epoch.  Sau mỗi epoch, ATE~F1 trên tập
phát triển được tính; checkpoint với F1 cao nhất được lưu làm kết quả
cuối cùng.

% ---------------------------------------------------------------------
\subsubsection{Hậu xử lý kết quả sinh}
\label{subsubsec:t5ate_postprocess}
% ---------------------------------------------------------------------

Chuỗi đầu ra được parse bằng biểu thức chính quy một nhóm:

\begin{equation}
  \texttt{\_ASPECT\_PATTERN} =
    \texttt{\textbackslash(([{\caret})]*)\textbackslash)}
  \label{eq:ate_regex}
\end{equation}

\noindent để trích xuất nội dung bên trong mỗi cặp ngoặc tròn thành
một aspect term.

% =====================================================================
\subsection{Chuẩn hóa aspect term bằng Levenshtein}
\label{subsubsec:levenshtein}
% =====================================================================

Cả hai mô hình đều áp dụng cùng một bước chuẩn hóa n-gram sau khi
parse, nhằm đảm bảo aspect term kết quả \textit{thực sự xuất hiện
trong câu gốc}:

\begin{enumerate}
  \item \textbf{Xây dựng từ điển n-gram:}  Từ câu gốc $\mathbf{x}$,
    tập ứng viên $\mathcal{V}$ gồm tất cả chuỗi con liên tục theo từ
    (unigram, bigram, trigram, \dots).

  \item \textbf{Ánh xạ về n-gram gần nhất:}  Nếu aspect term dự đoán
    $\hat{a}$ chưa thuộc $\mathcal{V}$, thay bằng:
    \begin{equation}
      a^* = \arg\min_{v \in \mathcal{V}}\;
        \mathrm{Lev}(\hat{a},\, v)
      \label{eq:levenshtein}
    \end{equation}
    \noindent trong đó $\mathrm{Lev}(\cdot, \cdot)$ là khoảng cách
    chỉnh sửa Levenshtein.  Nếu $\hat{a} \in \mathcal{V}$, giữ
    nguyên.
\end{enumerate}

% =====================================================================
\subsection{Độ đo đánh giá}
\label{subsec:ate_eval}
% =====================================================================

Cả hai phương pháp được đánh giá theo \textbf{ATE~F1}: một dự đoán
được tính là True Positive khi chuỗi aspect term dự đoán khớp chính
xác (\textit{exact match}) với một aspect term trong nhãn vàng của
câu.

Precision, Recall và F1 được tính theo công thức chuẩn:

\begin{equation}
  P = \frac{\mathrm{TP}}{\mathrm{TP} + \mathrm{FP}}, \quad
  R = \frac{\mathrm{TP}}{\mathrm{TP} + \mathrm{FN}}, \quad
  F_1 = \frac{2PR}{P + R}
  \label{eq:ate_prf}
\end{equation}

\noindent trong đó TP, FP, FN được tổng hợp (\textit{micro}) qua toàn
bộ câu của tập kiểm thử trước khi tính tỉ lệ — nhất quán với giao
thức đánh giá ATE trong các nghiên cứu ABSA~\cite{zhang2021gas,
luo2019doer}.
 trong main.tex
% =====================================================================

% =====================================================================
\section{Phương pháp T5 sinh cho Trích xuất Aspect Term}
\label{sec:ate}
% =====================================================================

Luận văn sử dụng hai cách tiếp cận dựa trên T5 để trích xuất aspect
term (ATE), đều được đánh giá theo cùng một độ đo ATE~F1:

\begin{enumerate}
  \item \textbf{GAS} (\textit{Generative Aspect Sentiment}) —
    mô hình seq2seq sinh đồng thời ba nhãn
    $(a_i, c_i, s_i)$ trong một lần giải mã; khi dùng cho ATE, chỉ
    trường $a_i$ được giữ lại
    (Mục~\ref{subsec:gas_full}).

  \item \textbf{T5-ATE} — mô hình seq2seq huấn luyện chuyên biệt cho
    ATE: chuỗi đích chỉ chứa aspect term, không có category hay
    sentiment (Mục~\ref{subsec:t5ate}).
\end{enumerate}

% =====================================================================
\subsection{Mô hình GAS (Generative Aspect Sentiment)}
\label{subsec:gas_full}
% =====================================================================

GAS~\cite{zhang2021gas} xây dựng bài toán ABSA như một bài toán dịch
chuỗi, với mô hình nền là
\texttt{T5ForConditionalGeneration}~\cite{raffel2020t5} phiên bản
\texttt{t5-base}.  Với mỗi câu đầu vào $\mathbf{x}$, mô hình được
huấn luyện để sinh ra chuỗi đích $\mathbf{y}$ theo cú pháp cố định:

\begin{equation}
  \mathbf{y} = \texttt{(}a_1,\, c_1,\, s_1\texttt{);\;
               (}a_2,\, c_2,\, s_2\texttt{);\;} \dots
  \label{eq:gas_format}
\end{equation}

\noindent trong đó $a_i$ là cụm từ khía cạnh (\textit{aspect term}),
$c_i \in \{\text{SERVICE, FACILITY, AMENITY, EXPERIENCE, BRANDING,
LOYALTY}\}$ là danh mục khía cạnh, và
$s_i \in \{\text{positive, negative, neutral}\}$ là cực tính cảm xúc.
Nếu câu không chứa khía cạnh nào, mô hình sinh ra chuỗi
\texttt{"none"}.

% ---------------------------------------------------------------------
\subsubsection{Tiền xử lý dữ liệu}
\label{subsubsec:gas_data}
% ---------------------------------------------------------------------

Dữ liệu đầu vào được đọc từ định dạng \texttt{.apc} bốn dòng lặp lại
(câu — aspect term — aspect category — sentiment).
Vì định dạng này lưu một dòng cho mỗi cặp \textit{(câu, aspect)},
các mẫu cùng thuộc một câu được nhóm lại thành một
\textit{record} duy nhất.  Chuỗi đích tương ứng được xây dựng bằng
cách nối tất cả bộ ba của câu đó theo cú pháp~\eqref{eq:gas_format}.

Mỗi record có dạng:

\begin{equation}
  r = \bigl(\underbrace{x}_{\text{input\_text}},\;
            \underbrace{y}_{\text{target\_text}},\;
            \underbrace{\{a_i\}}_{\text{aspects}},\;
            \underbrace{\{(a_i,\,c_i,\,s_i)\}}_{\text{gold\_triples}}
      \bigr)
  \label{eq:gas_record}
\end{equation}

Tokenization được thực hiện bởi \texttt{T5Tokenizer} với
\texttt{max\_length}$= 128$ cho cả đầu vào lẫn đầu ra.  Các vị trí
padding trong chuỗi nhãn được thay bằng $-100$ để hàm mất mát
cross-entropy bỏ qua.

% ---------------------------------------------------------------------
\subsubsection{Huấn luyện}
\label{subsubsec:gas_train}
% ---------------------------------------------------------------------

Mô hình được tinh chỉnh bằng hàm mất mát cross-entropy seq2seq chuẩn
(T5 tự tính theo cơ chế \textit{teacher forcing}):

\begin{equation}
  \mathcal{L}_{\text{GAS}} =
    -\sum_{t=1}^{T} \log P\!\left(y_t \mid y_{<t},\, \mathbf{x};\,
    \Theta\right)
  \label{eq:gas_loss}
\end{equation}

\noindent với $T$ là độ dài chuỗi đích và $\Theta$ là toàn bộ tham số
mô hình.  Quá trình tối ưu sử dụng AdamW với learning rate
$3 \times 10^{-4}$, kết hợp \textit{linear warmup} trong 10\% tổng
số bước.  Kỹ thuật mixed precision (AMP) được bật khi có GPU CUDA.

Early stopping theo dõi \textbf{joint F1} trên tập phát triển
(patience$= 4$ epoch): một dự đoán được tính là True Positive khi
\textit{cả ba} trường $(a_i, c_i, s_i)$ đều khớp chính xác với nhãn
vàng tương ứng.

% ---------------------------------------------------------------------
\subsubsection{Hậu xử lý và trích xuất aspect term}
\label{subsubsec:gas_postprocess}
% ---------------------------------------------------------------------

Chuỗi văn bản do T5 sinh ra được phân tích cú pháp bằng biểu thức
chính quy ba nhóm:

\begin{equation}
  \texttt{\_GAS\_RE} =
    \texttt{\textbackslash(\textbackslash s*(.+?),\textbackslash s*%
    (.+?),\textbackslash s*(positive|negative|neutral)\textbackslash s*%
    \textbackslash)}
  \label{eq:gas_regex}
\end{equation}

\noindent Group~1, 2, 3 lần lượt trả về $a_i$, $c_i$, $s_i$.  Để
sử dụng GAS như một baseline ATE, chỉ Group~1 ($a_i$) được giữ lại;
$c_i$ và $s_i$ bị bỏ qua ở bước đánh giá.  Nếu chuỗi sinh là
\texttt{"none"} hoặc không khớp mẫu nào, hàm trả về danh sách rỗng.

Sau khi parse, bước chuẩn hóa n-gram Levenshtein được áp dụng cho
$a_i$ (mục~\ref{subsubsec:levenshtein}); $c_i$ và $s_i$ giữ nguyên.

% =====================================================================
\subsection{Mô hình T5-ATE}
\label{subsec:t5ate}
% =====================================================================

Khác với GAS, T5-ATE được huấn luyện chuyên biệt cho ATE: chuỗi đích
chỉ chứa aspect term, không có category hay sentiment.  Điều này rút
ngắn không gian sinh và giảm áp lực cho decoder khi chỉ cần học một
nhãn.

% ---------------------------------------------------------------------
\subsubsection{Kiến trúc và định dạng đầu ra}
\label{subsubsec:t5ate_arch}
% ---------------------------------------------------------------------

Mô hình nền cũng là \texttt{T5ForConditionalGeneration} phiên bản
\texttt{t5-base}.  Với mỗi câu đầu vào $\mathbf{x}$, mô hình sinh ra
chuỗi đích chỉ chứa aspect term:

\begin{equation}
  \mathbf{y} = \texttt{(}a_1\texttt{);\;(}a_2\texttt{);\;} \dots
  \label{eq:ate_format}
\end{equation}

\noindent Nếu câu không chứa khía cạnh nào, mô hình sinh ra chuỗi
\texttt{"none"}.  Quá trình giải mã sử dụng \textit{beam search} với
chiều rộng tia $k = 4$.

% ---------------------------------------------------------------------
\subsubsection{Tiền xử lý dữ liệu}
\label{subsubsec:t5ate_data}
% ---------------------------------------------------------------------

Dữ liệu được đọc từ cùng định dạng \texttt{.apc} nhưng các trường
category và sentiment bị bỏ qua hoàn toàn.  Các mẫu cùng thuộc một
câu được nhóm lại; chuỗi đích được xây dựng theo
cú pháp~\eqref{eq:ate_format}.  Mỗi record có dạng:

\begin{equation}
  r = \bigl(\underbrace{x}_{\text{input\_text}},\;
            \underbrace{y}_{\text{target\_text}},\;
            \underbrace{\{a_i\}}_{\text{aspects}}
      \bigr)
  \label{eq:ate_record}
\end{equation}

Tokenization sử dụng \texttt{max\_length}$= 128$ cho đầu vào và
$= 64$ cho đầu ra — ngắn hơn GAS vì không cần mã hóa category và
sentiment.

% ---------------------------------------------------------------------
\subsubsection{Huấn luyện}
\label{subsubsec:t5ate_train}
% ---------------------------------------------------------------------

Hàm mất mát và thiết lập tối ưu giống GAS~\eqref{eq:gas_loss}:
AdamW với learning rate $3 \times 10^{-4}$ và linear warmup 10\%.
Điểm khác biệt: T5-ATE \textbf{không áp dụng early stopping} — mô
hình được huấn luyện đủ 20 epoch.  Sau mỗi epoch, ATE~F1 trên tập
phát triển được tính; checkpoint với F1 cao nhất được lưu làm kết quả
cuối cùng.

% ---------------------------------------------------------------------
\subsubsection{Hậu xử lý kết quả sinh}
\label{subsubsec:t5ate_postprocess}
% ---------------------------------------------------------------------

Chuỗi đầu ra được parse bằng biểu thức chính quy một nhóm:

\begin{equation}
  \texttt{\_ASPECT\_PATTERN} =
    \texttt{\textbackslash(([{\caret})]*)\textbackslash)}
  \label{eq:ate_regex}
\end{equation}

\noindent để trích xuất nội dung bên trong mỗi cặp ngoặc tròn thành
một aspect term.

% =====================================================================
\subsection{Chuẩn hóa aspect term bằng Levenshtein}
\label{subsubsec:levenshtein}
% =====================================================================

Cả hai mô hình đều áp dụng cùng một bước chuẩn hóa n-gram sau khi
parse, nhằm đảm bảo aspect term kết quả \textit{thực sự xuất hiện
trong câu gốc}:

\begin{enumerate}
  \item \textbf{Xây dựng từ điển n-gram:}  Từ câu gốc $\mathbf{x}$,
    tập ứng viên $\mathcal{V}$ gồm tất cả chuỗi con liên tục theo từ
    (unigram, bigram, trigram, \dots).

  \item \textbf{Ánh xạ về n-gram gần nhất:}  Nếu aspect term dự đoán
    $\hat{a}$ chưa thuộc $\mathcal{V}$, thay bằng:
    \begin{equation}
      a^* = \arg\min_{v \in \mathcal{V}}\;
        \mathrm{Lev}(\hat{a},\, v)
      \label{eq:levenshtein}
    \end{equation}
    \noindent trong đó $\mathrm{Lev}(\cdot, \cdot)$ là khoảng cách
    chỉnh sửa Levenshtein.  Nếu $\hat{a} \in \mathcal{V}$, giữ
    nguyên.
\end{enumerate}

% =====================================================================
\subsection{Độ đo đánh giá}
\label{subsec:ate_eval}
% =====================================================================

Cả hai phương pháp được đánh giá theo \textbf{ATE~F1}: một dự đoán
được tính là True Positive khi chuỗi aspect term dự đoán khớp chính
xác (\textit{exact match}) với một aspect term trong nhãn vàng của
câu.

Precision, Recall và F1 được tính theo công thức chuẩn:

\begin{equation}
  P = \frac{\mathrm{TP}}{\mathrm{TP} + \mathrm{FP}}, \quad
  R = \frac{\mathrm{TP}}{\mathrm{TP} + \mathrm{FN}}, \quad
  F_1 = \frac{2PR}{P + R}
  \label{eq:ate_prf}
\end{equation}

\noindent trong đó TP, FP, FN được tổng hợp (\textit{micro}) qua toàn
bộ câu của tập kiểm thử trước khi tính tỉ lệ — nhất quán với giao
thức đánh giá ATE trong các nghiên cứu ABSA~\cite{zhang2021gas,
luo2019doer}.
 trong main.tex
% =====================================================================

% =====================================================================
\section{Phương pháp T5 sinh cho Trích xuất Aspect Term}
\label{sec:ate}
% =====================================================================

Luận văn sử dụng hai cách tiếp cận dựa trên T5 để trích xuất aspect
term (ATE), đều được đánh giá theo cùng một độ đo ATE~F1:

\begin{enumerate}
  \item \textbf{GAS} (\textit{Generative Aspect Sentiment}) —
    mô hình seq2seq sinh đồng thời ba nhãn
    $(a_i, c_i, s_i)$ trong một lần giải mã; khi dùng cho ATE, chỉ
    trường $a_i$ được giữ lại
    (Mục~\ref{subsec:gas_full}).

  \item \textbf{T5-ATE} — mô hình seq2seq huấn luyện chuyên biệt cho
    ATE: chuỗi đích chỉ chứa aspect term, không có category hay
    sentiment (Mục~\ref{subsec:t5ate}).
\end{enumerate}

% =====================================================================
\subsection{Mô hình GAS (Generative Aspect Sentiment)}
\label{subsec:gas_full}
% =====================================================================

GAS~\cite{zhang2021gas} xây dựng bài toán ABSA như một bài toán dịch
chuỗi, với mô hình nền là
\texttt{T5ForConditionalGeneration}~\cite{raffel2020t5} phiên bản
\texttt{t5-base}.  Với mỗi câu đầu vào $\mathbf{x}$, mô hình được
huấn luyện để sinh ra chuỗi đích $\mathbf{y}$ theo cú pháp cố định:

\begin{equation}
  \mathbf{y} = \texttt{(}a_1,\, c_1,\, s_1\texttt{);\;
               (}a_2,\, c_2,\, s_2\texttt{);\;} \dots
  \label{eq:gas_format}
\end{equation}

\noindent trong đó $a_i$ là cụm từ khía cạnh (\textit{aspect term}),
$c_i \in \{\text{SERVICE, FACILITY, AMENITY, EXPERIENCE, BRANDING,
LOYALTY}\}$ là danh mục khía cạnh, và
$s_i \in \{\text{positive, negative, neutral}\}$ là cực tính cảm xúc.
Nếu câu không chứa khía cạnh nào, mô hình sinh ra chuỗi
\texttt{"none"}.

% ---------------------------------------------------------------------
\subsubsection{Tiền xử lý dữ liệu}
\label{subsubsec:gas_data}
% ---------------------------------------------------------------------

Dữ liệu đầu vào được đọc từ định dạng \texttt{.apc} bốn dòng lặp lại
(câu — aspect term — aspect category — sentiment).
Vì định dạng này lưu một dòng cho mỗi cặp \textit{(câu, aspect)},
các mẫu cùng thuộc một câu được nhóm lại thành một
\textit{record} duy nhất.  Chuỗi đích tương ứng được xây dựng bằng
cách nối tất cả bộ ba của câu đó theo cú pháp~\eqref{eq:gas_format}.

Mỗi record có dạng:

\begin{equation}
  r = \bigl(\underbrace{x}_{\text{input\_text}},\;
            \underbrace{y}_{\text{target\_text}},\;
            \underbrace{\{a_i\}}_{\text{aspects}},\;
            \underbrace{\{(a_i,\,c_i,\,s_i)\}}_{\text{gold\_triples}}
      \bigr)
  \label{eq:gas_record}
\end{equation}

Tokenization được thực hiện bởi \texttt{T5Tokenizer} với
\texttt{max\_length}$= 128$ cho cả đầu vào lẫn đầu ra.  Các vị trí
padding trong chuỗi nhãn được thay bằng $-100$ để hàm mất mát
cross-entropy bỏ qua.

% ---------------------------------------------------------------------
\subsubsection{Huấn luyện}
\label{subsubsec:gas_train}
% ---------------------------------------------------------------------

Mô hình được tinh chỉnh bằng hàm mất mát cross-entropy seq2seq chuẩn
(T5 tự tính theo cơ chế \textit{teacher forcing}):

\begin{equation}
  \mathcal{L}_{\text{GAS}} =
    -\sum_{t=1}^{T} \log P\!\left(y_t \mid y_{<t},\, \mathbf{x};\,
    \Theta\right)
  \label{eq:gas_loss}
\end{equation}

\noindent với $T$ là độ dài chuỗi đích và $\Theta$ là toàn bộ tham số
mô hình.  Quá trình tối ưu sử dụng AdamW với learning rate
$3 \times 10^{-4}$, kết hợp \textit{linear warmup} trong 10\% tổng
số bước.  Kỹ thuật mixed precision (AMP) được bật khi có GPU CUDA.

Early stopping theo dõi \textbf{joint F1} trên tập phát triển
(patience$= 4$ epoch): một dự đoán được tính là True Positive khi
\textit{cả ba} trường $(a_i, c_i, s_i)$ đều khớp chính xác với nhãn
vàng tương ứng.

% ---------------------------------------------------------------------
\subsubsection{Hậu xử lý và trích xuất aspect term}
\label{subsubsec:gas_postprocess}
% ---------------------------------------------------------------------

Chuỗi văn bản do T5 sinh ra được phân tích cú pháp bằng biểu thức
chính quy ba nhóm:

\begin{equation}
  \texttt{\_GAS\_RE} =
    \texttt{\textbackslash(\textbackslash s*(.+?),\textbackslash s*%
    (.+?),\textbackslash s*(positive|negative|neutral)\textbackslash s*%
    \textbackslash)}
  \label{eq:gas_regex}
\end{equation}

\noindent Group~1, 2, 3 lần lượt trả về $a_i$, $c_i$, $s_i$.  Để
sử dụng GAS như một baseline ATE, chỉ Group~1 ($a_i$) được giữ lại;
$c_i$ và $s_i$ bị bỏ qua ở bước đánh giá.  Nếu chuỗi sinh là
\texttt{"none"} hoặc không khớp mẫu nào, hàm trả về danh sách rỗng.

Sau khi parse, bước chuẩn hóa n-gram Levenshtein được áp dụng cho
$a_i$ (mục~\ref{subsubsec:levenshtein}); $c_i$ và $s_i$ giữ nguyên.

% =====================================================================
\subsection{Mô hình T5-ATE}
\label{subsec:t5ate}
% =====================================================================

Khác với GAS, T5-ATE được huấn luyện chuyên biệt cho ATE: chuỗi đích
chỉ chứa aspect term, không có category hay sentiment.  Điều này rút
ngắn không gian sinh và giảm áp lực cho decoder khi chỉ cần học một
nhãn.

% ---------------------------------------------------------------------
\subsubsection{Kiến trúc và định dạng đầu ra}
\label{subsubsec:t5ate_arch}
% ---------------------------------------------------------------------

Mô hình nền cũng là \texttt{T5ForConditionalGeneration} phiên bản
\texttt{t5-base}.  Với mỗi câu đầu vào $\mathbf{x}$, mô hình sinh ra
chuỗi đích chỉ chứa aspect term:

\begin{equation}
  \mathbf{y} = \texttt{(}a_1\texttt{);\;(}a_2\texttt{);\;} \dots
  \label{eq:ate_format}
\end{equation}

\noindent Nếu câu không chứa khía cạnh nào, mô hình sinh ra chuỗi
\texttt{"none"}.  Quá trình giải mã sử dụng \textit{beam search} với
chiều rộng tia $k = 4$.

% ---------------------------------------------------------------------
\subsubsection{Tiền xử lý dữ liệu}
\label{subsubsec:t5ate_data}
% ---------------------------------------------------------------------

Dữ liệu được đọc từ cùng định dạng \texttt{.apc} nhưng các trường
category và sentiment bị bỏ qua hoàn toàn.  Các mẫu cùng thuộc một
câu được nhóm lại; chuỗi đích được xây dựng theo
cú pháp~\eqref{eq:ate_format}.  Mỗi record có dạng:

\begin{equation}
  r = \bigl(\underbrace{x}_{\text{input\_text}},\;
            \underbrace{y}_{\text{target\_text}},\;
            \underbrace{\{a_i\}}_{\text{aspects}}
      \bigr)
  \label{eq:ate_record}
\end{equation}

Tokenization sử dụng \texttt{max\_length}$= 128$ cho đầu vào và
$= 64$ cho đầu ra — ngắn hơn GAS vì không cần mã hóa category và
sentiment.

% ---------------------------------------------------------------------
\subsubsection{Huấn luyện}
\label{subsubsec:t5ate_train}
% ---------------------------------------------------------------------

Hàm mất mát và thiết lập tối ưu giống GAS~\eqref{eq:gas_loss}:
AdamW với learning rate $3 \times 10^{-4}$ và linear warmup 10\%.
Điểm khác biệt: T5-ATE \textbf{không áp dụng early stopping} — mô
hình được huấn luyện đủ 20 epoch.  Sau mỗi epoch, ATE~F1 trên tập
phát triển được tính; checkpoint với F1 cao nhất được lưu làm kết quả
cuối cùng.

% ---------------------------------------------------------------------
\subsubsection{Hậu xử lý kết quả sinh}
\label{subsubsec:t5ate_postprocess}
% ---------------------------------------------------------------------

Chuỗi đầu ra được parse bằng biểu thức chính quy một nhóm:

\begin{equation}
  \texttt{\_ASPECT\_PATTERN} =
    \texttt{\textbackslash(([{\caret})]*)\textbackslash)}
  \label{eq:ate_regex}
\end{equation}

\noindent để trích xuất nội dung bên trong mỗi cặp ngoặc tròn thành
một aspect term.

% =====================================================================
\subsection{Chuẩn hóa aspect term bằng Levenshtein}
\label{subsubsec:levenshtein}
% =====================================================================

Cả hai mô hình đều áp dụng cùng một bước chuẩn hóa n-gram sau khi
parse, nhằm đảm bảo aspect term kết quả \textit{thực sự xuất hiện
trong câu gốc}:

\begin{enumerate}
  \item \textbf{Xây dựng từ điển n-gram:}  Từ câu gốc $\mathbf{x}$,
    tập ứng viên $\mathcal{V}$ gồm tất cả chuỗi con liên tục theo từ
    (unigram, bigram, trigram, \dots).

  \item \textbf{Ánh xạ về n-gram gần nhất:}  Nếu aspect term dự đoán
    $\hat{a}$ chưa thuộc $\mathcal{V}$, thay bằng:
    \begin{equation}
      a^* = \arg\min_{v \in \mathcal{V}}\;
        \mathrm{Lev}(\hat{a},\, v)
      \label{eq:levenshtein}
    \end{equation}
    \noindent trong đó $\mathrm{Lev}(\cdot, \cdot)$ là khoảng cách
    chỉnh sửa Levenshtein.  Nếu $\hat{a} \in \mathcal{V}$, giữ
    nguyên.
\end{enumerate}

% =====================================================================
\subsection{Độ đo đánh giá}
\label{subsec:ate_eval}
% =====================================================================

Cả hai phương pháp được đánh giá theo \textbf{ATE~F1}: một dự đoán
được tính là True Positive khi chuỗi aspect term dự đoán khớp chính
xác (\textit{exact match}) với một aspect term trong nhãn vàng của
câu.

Precision, Recall và F1 được tính theo công thức chuẩn:

\begin{equation}
  P = \frac{\mathrm{TP}}{\mathrm{TP} + \mathrm{FP}}, \quad
  R = \frac{\mathrm{TP}}{\mathrm{TP} + \mathrm{FN}}, \quad
  F_1 = \frac{2PR}{P + R}
  \label{eq:ate_prf}
\end{equation}

\noindent trong đó TP, FP, FN được tổng hợp (\textit{micro}) qua toàn
bộ câu của tập kiểm thử trước khi tính tỉ lệ — nhất quán với giao
thức đánh giá ATE trong các nghiên cứu ABSA~\cite{zhang2021gas,
luo2019doer}.
 trong main.tex
% =====================================================================

% =====================================================================
\section{Phương pháp T5 sinh cho Trích xuất Aspect Term}
\label{sec:ate}
% =====================================================================

Luận văn sử dụng hai cách tiếp cận dựa trên T5 để trích xuất aspect
term (ATE), đều được đánh giá theo cùng một độ đo ATE~F1:

\begin{enumerate}
  \item \textbf{GAS} (\textit{Generative Aspect Sentiment}) —
    mô hình seq2seq sinh đồng thời ba nhãn
    $(a_i, c_i, s_i)$ trong một lần giải mã; khi dùng cho ATE, chỉ
    trường $a_i$ được giữ lại
    (Mục~\ref{subsec:gas_full}).

  \item \textbf{T5-ATE} — mô hình seq2seq huấn luyện chuyên biệt cho
    ATE: chuỗi đích chỉ chứa aspect term, không có category hay
    sentiment (Mục~\ref{subsec:t5ate}).
\end{enumerate}

% =====================================================================
\subsection{Mô hình GAS (Generative Aspect Sentiment)}
\label{subsec:gas_full}
% =====================================================================

GAS~\cite{zhang2021gas} xây dựng bài toán ABSA như một bài toán dịch
chuỗi, với mô hình nền là
\texttt{T5ForConditionalGeneration}~\cite{raffel2020t5} phiên bản
\texttt{t5-base}.  Với mỗi câu đầu vào $\mathbf{x}$, mô hình được
huấn luyện để sinh ra chuỗi đích $\mathbf{y}$ theo cú pháp cố định:

\begin{equation}
  \mathbf{y} = \texttt{(}a_1,\, c_1,\, s_1\texttt{);\;
               (}a_2,\, c_2,\, s_2\texttt{);\;} \dots
  \label{eq:gas_format}
\end{equation}

\noindent trong đó $a_i$ là cụm từ khía cạnh (\textit{aspect term}),
$c_i \in \{\text{SERVICE, FACILITY, AMENITY, EXPERIENCE, BRANDING,
LOYALTY}\}$ là danh mục khía cạnh, và
$s_i \in \{\text{positive, negative, neutral}\}$ là cực tính cảm xúc.
Nếu câu không chứa khía cạnh nào, mô hình sinh ra chuỗi
\texttt{"none"}.

% ---------------------------------------------------------------------
\subsubsection{Tiền xử lý dữ liệu}
\label{subsubsec:gas_data}
% ---------------------------------------------------------------------

Dữ liệu đầu vào được đọc từ định dạng \texttt{.apc} bốn dòng lặp lại
(câu — aspect term — aspect category — sentiment).
Vì định dạng này lưu một dòng cho mỗi cặp \textit{(câu, aspect)},
các mẫu cùng thuộc một câu được nhóm lại thành một
\textit{record} duy nhất.  Chuỗi đích tương ứng được xây dựng bằng
cách nối tất cả bộ ba của câu đó theo cú pháp~\eqref{eq:gas_format}.

Mỗi record có dạng:

\begin{equation}
  r = \bigl(\underbrace{x}_{\text{input\_text}},\;
            \underbrace{y}_{\text{target\_text}},\;
            \underbrace{\{a_i\}}_{\text{aspects}},\;
            \underbrace{\{(a_i,\,c_i,\,s_i)\}}_{\text{gold\_triples}}
      \bigr)
  \label{eq:gas_record}
\end{equation}

Tokenization được thực hiện bởi \texttt{T5Tokenizer} với
\texttt{max\_length}$= 128$ cho cả đầu vào lẫn đầu ra.  Các vị trí
padding trong chuỗi nhãn được thay bằng $-100$ để hàm mất mát
cross-entropy bỏ qua.

% ---------------------------------------------------------------------
\subsubsection{Huấn luyện}
\label{subsubsec:gas_train}
% ---------------------------------------------------------------------

Mô hình được tinh chỉnh bằng hàm mất mát cross-entropy seq2seq chuẩn
(T5 tự tính theo cơ chế \textit{teacher forcing}):

\begin{equation}
  \mathcal{L}_{\text{GAS}} =
    -\sum_{t=1}^{T} \log P\!\left(y_t \mid y_{<t},\, \mathbf{x};\,
    \Theta\right)
  \label{eq:gas_loss}
\end{equation}

\noindent với $T$ là độ dài chuỗi đích và $\Theta$ là toàn bộ tham số
mô hình.  Quá trình tối ưu sử dụng AdamW với learning rate
$3 \times 10^{-4}$, kết hợp \textit{linear warmup} trong 10\% tổng
số bước.  Kỹ thuật mixed precision (AMP) được bật khi có GPU CUDA.

Early stopping theo dõi \textbf{joint F1} trên tập phát triển
(patience$= 4$ epoch): một dự đoán được tính là True Positive khi
\textit{cả ba} trường $(a_i, c_i, s_i)$ đều khớp chính xác với nhãn
vàng tương ứng.

% ---------------------------------------------------------------------
\subsubsection{Hậu xử lý và trích xuất aspect term}
\label{subsubsec:gas_postprocess}
% ---------------------------------------------------------------------

Chuỗi văn bản do T5 sinh ra được phân tích cú pháp bằng biểu thức
chính quy ba nhóm:

\begin{equation}
  \texttt{\_GAS\_RE} =
    \texttt{\textbackslash(\textbackslash s*(.+?),\textbackslash s*%
    (.+?),\textbackslash s*(positive|negative|neutral)\textbackslash s*%
    \textbackslash)}
  \label{eq:gas_regex}
\end{equation}

\noindent Group~1, 2, 3 lần lượt trả về $a_i$, $c_i$, $s_i$.  Để
sử dụng GAS như một baseline ATE, chỉ Group~1 ($a_i$) được giữ lại;
$c_i$ và $s_i$ bị bỏ qua ở bước đánh giá.  Nếu chuỗi sinh là
\texttt{"none"} hoặc không khớp mẫu nào, hàm trả về danh sách rỗng.

Sau khi parse, bước chuẩn hóa n-gram Levenshtein được áp dụng cho
$a_i$ (mục~\ref{subsubsec:levenshtein}); $c_i$ và $s_i$ giữ nguyên.

% =====================================================================
\subsection{Mô hình T5-ATE}
\label{subsec:t5ate}
% =====================================================================

Khác với GAS, T5-ATE được huấn luyện chuyên biệt cho ATE: chuỗi đích
chỉ chứa aspect term, không có category hay sentiment.  Điều này rút
ngắn không gian sinh và giảm áp lực cho decoder khi chỉ cần học một
nhãn.

% ---------------------------------------------------------------------
\subsubsection{Kiến trúc và định dạng đầu ra}
\label{subsubsec:t5ate_arch}
% ---------------------------------------------------------------------

Mô hình nền cũng là \texttt{T5ForConditionalGeneration} phiên bản
\texttt{t5-base}.  Với mỗi câu đầu vào $\mathbf{x}$, mô hình sinh ra
chuỗi đích chỉ chứa aspect term:

\begin{equation}
  \mathbf{y} = \texttt{(}a_1\texttt{);\;(}a_2\texttt{);\;} \dots
  \label{eq:ate_format}
\end{equation}

\noindent Nếu câu không chứa khía cạnh nào, mô hình sinh ra chuỗi
\texttt{"none"}.  Quá trình giải mã sử dụng \textit{beam search} với
chiều rộng tia $k = 4$.

% ---------------------------------------------------------------------
\subsubsection{Tiền xử lý dữ liệu}
\label{subsubsec:t5ate_data}
% ---------------------------------------------------------------------

Dữ liệu được đọc từ cùng định dạng \texttt{.apc} nhưng các trường
category và sentiment bị bỏ qua hoàn toàn.  Các mẫu cùng thuộc một
câu được nhóm lại; chuỗi đích được xây dựng theo
cú pháp~\eqref{eq:ate_format}.  Mỗi record có dạng:

\begin{equation}
  r = \bigl(\underbrace{x}_{\text{input\_text}},\;
            \underbrace{y}_{\text{target\_text}},\;
            \underbrace{\{a_i\}}_{\text{aspects}}
      \bigr)
  \label{eq:ate_record}
\end{equation}

Tokenization sử dụng \texttt{max\_length}$= 128$ cho đầu vào và
$= 64$ cho đầu ra — ngắn hơn GAS vì không cần mã hóa category và
sentiment.

% ---------------------------------------------------------------------
\subsubsection{Huấn luyện}
\label{subsubsec:t5ate_train}
% ---------------------------------------------------------------------

Hàm mất mát và thiết lập tối ưu giống GAS~\eqref{eq:gas_loss}:
AdamW với learning rate $3 \times 10^{-4}$ và linear warmup 10\%.
Điểm khác biệt: T5-ATE \textbf{không áp dụng early stopping} — mô
hình được huấn luyện đủ 20 epoch.  Sau mỗi epoch, ATE~F1 trên tập
phát triển được tính; checkpoint với F1 cao nhất được lưu làm kết quả
cuối cùng.

% ---------------------------------------------------------------------
\subsubsection{Hậu xử lý kết quả sinh}
\label{subsubsec:t5ate_postprocess}
% ---------------------------------------------------------------------

Chuỗi đầu ra được parse bằng biểu thức chính quy một nhóm:

\begin{equation}
  \texttt{\_ASPECT\_PATTERN} =
    \texttt{\textbackslash(([{\caret})]*)\textbackslash)}
  \label{eq:ate_regex}
\end{equation}

\noindent để trích xuất nội dung bên trong mỗi cặp ngoặc tròn thành
một aspect term.

% =====================================================================
\subsection{Chuẩn hóa aspect term bằng Levenshtein}
\label{subsubsec:levenshtein}
% =====================================================================

Cả hai mô hình đều áp dụng cùng một bước chuẩn hóa n-gram sau khi
parse, nhằm đảm bảo aspect term kết quả \textit{thực sự xuất hiện
trong câu gốc}:

\begin{enumerate}
  \item \textbf{Xây dựng từ điển n-gram:}  Từ câu gốc $\mathbf{x}$,
    tập ứng viên $\mathcal{V}$ gồm tất cả chuỗi con liên tục theo từ
    (unigram, bigram, trigram, \dots).

  \item \textbf{Ánh xạ về n-gram gần nhất:}  Nếu aspect term dự đoán
    $\hat{a}$ chưa thuộc $\mathcal{V}$, thay bằng:
    \begin{equation}
      a^* = \arg\min_{v \in \mathcal{V}}\;
        \mathrm{Lev}(\hat{a},\, v)
      \label{eq:levenshtein}
    \end{equation}
    \noindent trong đó $\mathrm{Lev}(\cdot, \cdot)$ là khoảng cách
    chỉnh sửa Levenshtein.  Nếu $\hat{a} \in \mathcal{V}$, giữ
    nguyên.
\end{enumerate}

% =====================================================================
\subsection{Độ đo đánh giá}
\label{subsec:ate_eval}
% =====================================================================

Cả hai phương pháp được đánh giá theo \textbf{ATE~F1}: một dự đoán
được tính là True Positive khi chuỗi aspect term dự đoán khớp chính
xác (\textit{exact match}) với một aspect term trong nhãn vàng của
câu.

Precision, Recall và F1 được tính theo công thức chuẩn:

\begin{equation}
  P = \frac{\mathrm{TP}}{\mathrm{TP} + \mathrm{FP}}, \quad
  R = \frac{\mathrm{TP}}{\mathrm{TP} + \mathrm{FN}}, \quad
  F_1 = \frac{2PR}{P + R}
  \label{eq:ate_prf}
\end{equation}

\noindent trong đó TP, FP, FN được tổng hợp (\textit{micro}) qua toàn
bộ câu của tập kiểm thử trước khi tính tỉ lệ — nhất quán với giao
thức đánh giá ATE trong các nghiên cứu ABSA~\cite{zhang2021gas,
luo2019doer}.
